\section{In-Context Learning}
One of the most surprising capabilities of modern LLMs is in-context learning \cite{zhao2025incontextlearningsufficientinstruction}, which emerges as a by-product of scale and the Transformer architecture. Unlike traditional machine learning systems, where models must be explicitly retrained for new tasks, LLMs can learn to perform novel tasks on the fly by conditioning on examples provided in the input prompt. For instance, if a prompt contains several examples of movie reviews labeled with "positive sentiment" or "negative sentiment," the model can infer the classification pattern and correctly label new reviews without any parameter updates. This phenomenon demonstrates that the model has implicitly acquired the ability to generalize from patterns in its input sequence, effectively treating the prompt as a temporary training set. In-context learning is a key enabler of “few-shot” and “zero-shot” performance, where models adapt to new domains or tasks with little or no additional supervision. It represents a fundamental shift in how AI systems can be deployed, offering adaptability without retraining and significantly expanding the practical versatility of LLMs \cite{han2023incontext, ibm2021incontext, googleresearch2023icl}.
In this thesis, this capability was massively exploited.